\documentclass[12pt]{article}
\usepackage[margin=1in]{geometry}
\usepackage[all]{xy}


\usepackage{amsmath,amsthm,amssymb,color,latexsym}
\usepackage{geometry}        
\geometry{letterpaper}  
\usepackage{boxedminipage}
\usepackage{graphicx}
\usepackage{hyperref}
\hypersetup{
    colorlinks=true,
    citecolor=blue,
    linkcolor=blue,
    filecolor=cyan,
    urlcolor=magenta,
}
\usepackage[capitalize,nameinlink]{cleveref}

\newtheorem{problem}{Problem}


\newenvironment{solution}[1][\it{Solution}]{\textbf{#1. } }{$\square$}
\newcommand{\pointsbox}[1]{%
  \colorbox{red!70}{\strut\textcolor{white}{\sffamily\bfseries\ \ #1\ \ }}%
}
\newcommand{\pagetag}[1]{%
  \colorbox{red!70}{\strut\textcolor{white}{\sffamily\bfseries\ \ #1\ \ }}%
}
\newcommand{\striptitle}[2]{%
  \noindent
  \begin{minipage}{0.07\linewidth}\centering \pagetag{#1}\end{minipage}%
  \hfill
  \begin{minipage}{0.89\linewidth}
    {\large\scshape #2}
  \end{minipage}\par\vspace{0.75em}
}
% --- Environments ---
\theoremstyle{plain}
\newtheorem{theorem}{Theorem}
\newtheorem{corollary}[theorem]{Corollary}
\theoremstyle{definition}
\newtheorem{definition}[theorem]{Definition}
\newtheorem{remark}[theorem]{Remark}

\newenvironment{boxedalgo}
  {\begin{center}\begin{boxedminipage}{1\linewidth}}
  {\end{boxedminipage}\end{center}}

\newenvironment{inlinealgo}[1]
  {\noindent\hrulefill\\\noindent{\bf {#1}}}
  {\hrulefill}


% --- Shortcuts ---
\newcommand{\R}{\mathbb{R}}
\newcommand{\inner}[2]{\left\langle #1,#2 \right\rangle}
\newcommand{\proj}{\mathrm{proj}}
\newcommand{\Id}{\mathbf{I}}
\newcommand{\ot}{\otimes}
\newcommand{\mc}{\mathcal}

\begin{document}
\noindent CS 578 Statistical Machine Learning (Fall 2026)\hfill Problem Set \# 2\\
\{\{Your name\}\} \hfill \textbf{Due: September 16th}

\noindent\hrulefill

\subsection*{Problem 1. Averaging argument.}
In the lecture, we have seen the following version of the ``No free lunch'' theorem:
\begin{theorem}\label{thm:NFL}
Consider binary classification with 0–1 loss and ${\cal X}$ infinite. For any $m>0$ and any learning algorithm $A$, there exists a probability distribution ${\cal D}$ such that
\begin{align*}
    {\cal R}_\star=0,\qquad \mathbb{E}_{S\sim {\cal D}^m}[\mathcal{R}(A(S))]\geq \frac{1}{4}.
\end{align*}
\end{theorem}
You will prove the following corollary using~\cref{thm:NFL}: 
\begin{corollary}
Let ${\cal X}$ be infinite and let ${\cal H}$ be the set of all functions from ${\cal X}$ to $\{0,1\}$. Then ${\cal H}$ is not PAC-learnable.
\end{corollary}
\noindent{\bfseries Solution.}\newline 
\   %%% <-- ERASE THIS LINE AND WRITE YOUR SOLUTION HERE
\newpage

\subsection*{Problem 2. ERM for learning closed intervals.}
Let ${\cal X}=\mathbb{R}$ and ${\cal H}$ be the set of indicator functions of closed intervals and the empty set:
\begin{align*}
    {\cal H}=\{\mathbf{1}_{a\leq x\leq b}:\forall a,b\in \mathbb{R}\}.
\end{align*}
Assume binary classification with $0$--$1$ loss and realizability. Prove that ${\cal H}$ can be PAC-learned using ERM with sample complexity
\begin{align*}
    O\left(\frac{\log (1/\delta)}{\epsilon}\right).
\end{align*}
You need to provide a description of the learning algorithm and prove its sample complexity and the PAC learning guarantee.

\noindent{\bfseries Solution.}\newline 
\   %%% <-- ERASE THIS LINE AND WRITE YOUR SOLUTION HERE
\newpage


\subsection*{Problem 3. VC dimensions.}
Given some finite domain set, ${\cal X}$, and a number $k\leq |{\cal X}|$, figure out the VC dimension of each of the following classes (and prove your claims):
\begin{enumerate}
    \item ${\cal H}^{\cal X}_{=k}=\left\{h\in \{0,1\}^{\cal X}:|\{x:h(x)=1\}|=k\right\}$. That is, the set of all functions that assign the value $1$ to exactly $k$ elements of ${\cal X}$.
    \item ${\cal H}_{\text{at-most-}k} = \left\{h\in \{0, 1\}^{\cal X}:|\{x:h(x)=1\}|\leq k~~\text{or}~~|\{x:h(x)=0\}|\leq k\right\}$.
\end{enumerate}
\end{document}
